\chapter{数据集构建与在线团雾合成方法}

\section{原始数据集选择与分析}

本文选用 UA-DETRAC 作为基础交通视频数据来源。该数据集具有以下优点：第一，数据来自固定道路监控摄像头，与本文研究的高速公路监控场景在视角上更接近；第二，包含连续视频帧和配套 XML 标注，便于开展序列级划分与检测监督构建；第三，车辆目标丰富，适合用作单类 \texttt{vehicle} 检测任务的数据基础。虽然 UA-DETRAC 并不是专门的雾天数据集，但其清晰天气样本丰富、结构规范，因此非常适合作为“清晰交通数据 + 在线造雾”路线中的基底数据。

根据当前工作区运行时统计结果，原始训练数据目录中共包含 60 个监控序列，测试目录中包含 40 个监控序列。按照当前代码中的随机种子、帧步长和序列级划分策略，运行时实际得到训练集 48 个序列、67790 个样本，验证集 12 个序列、16001 个样本。进一步统计可知，当前训练集包含 66811 个有标注样本、473719 个检测框，验证集包含 15274 个有标注样本、124562 个检测框。这些统计来自本文项目在当前环境中的真实运行，而不是静态文档估计值，因此更能反映论文撰写时系统的真实数据基础。

这组数据规模对于毕业设计和工程原型而言已经具备较高的可用性。一方面，车辆目标数量足够支撑单类检测训练；另一方面，固定监控视频序列的数据组织形式，十分适合研究“相同机位、不同气象退化下的联合感知问题”。然而，该数据集也存在局限：它缺少真实团雾标签，无法直接监督雾分类与浓度回归；此外，由于原始标注只覆盖车辆目标，因此天气任务必须通过合成方式构建监督信号。

\begin{figure}[htbp]
  \centering
  \includegraphics[width=0.95\textwidth]{fig_04_dataset_statistics.png}
  \caption{数据集构成统计图与统计表}
  \label{fig:dataset_statistics}
\end{figure}

\section{序列级训练集与验证集划分}

在交通监控视频场景中，若按图像帧随机划分训练集和验证集，会使同一摄像头、同一时段甚至相邻帧同时出现在训练集和验证集中，从而导致严重的数据泄漏问题。模型可能因为记住了背景结构、车道位置和局部目标外观，而在验证阶段获得虚高结果。因此，本文在 \texttt{src/utils.py} 中实现了按序列名排序、打乱并基于随机种子切分的 \texttt{split\_sequence\_names()} 逻辑，在 \texttt{MultiTaskDataset} 中严格按视频序列划分训练集和验证集。

这种划分方式有两个好处。其一，它保证了训练集与验证集之间不存在相邻帧泄漏，更符合视频任务的真实评估要求。其二，它使深度缓存预计算、数据审计脚本和训练主流程共享同一划分规则，从而避免“训练划分一套，预计算划分另一套”的工程混乱。当前系统中，\texttt{Config.TRAIN\_RATIO = 0.8}，随机种子为 42，因而每次重新构建数据集时都可得到稳定一致的 train/val 拆分结果。

\section{深度缓存预计算与管理}

由于本文在线造雾依赖场景深度信息，而训练集规模较大，若在每个 epoch 中为每一张图重复运行 MiDaS 推理，将显著增加训练时间并带来额外不稳定因素。因此，本文采用“训练前离线预计算 + 训练时直接读取缓存”的方式管理深度图。

具体实现为：首先调用 \texttt{DepthEstimator} 使用 MiDaS\_small 对每张图像估计相对深度，随后将深度图插值回原图尺寸、归一化到 \texttt{[0,1]} 区间，并统一处理为“0 表示近、1 表示远”的深度约定，再以 \texttt{.npy} 格式保存到 \texttt{outputs/Depth\_Cache/}。缓存文件命名规则为 \texttt{序列名\_图像名.npy}，例如某序列中的 \texttt{img00001.jpg} 对应缓存文件 \texttt{MVI\_xxxxx\_img00001.jpg.npy}。这种命名规则使图像、序列和深度缓存之间的映射关系简单直接，便于在 \texttt{\_\_getitem\_\_()} 中快速定位。

从当前运行情况来看，深度缓存覆盖已经完整。工作区中 \texttt{outputs/Depth\_Cache} 的 \texttt{.npy} 文件数量与原始训练图像总量一致，训练前 smoke test 也确认在抽查范围内训练集和验证集未发现缺失缓存。训练脚本还进一步实现了自动补算策略：即使用户没有手动执行深度预计算，只要发现 train/val 任一侧存在缺失缓存，也会在正式训练前自动补齐。这种设计显著增强了系统鲁棒性，降低了“训练半途因缺文件崩溃”的风险。

\section{多任务数据集设计}

本文的多任务数据集定义在 \texttt{src/data/dataset.py} 中，其核心目标并非输出离线雾图，而是为在线造雾和统一模型训练提供完整输入。每个样本返回四类信息：清晰图像张量、对应深度图张量、检测类别张量和检测框张量。其中，图像与深度图用于在线生成雾样本，检测类别与框则用于构造检测监督。

在具体实现中，数据集完成了以下几项关键工作：

\begin{enumerate}[leftmargin=2.5em]
\item 从原始序列目录扫描样本并建立 \code{(img_path, seq, img_name)} 三元组索引；
\item 预加载各序列 XML 标注，建立按 \code{sequence}、\code{frame_num} 与 \code{boxes} 组织的层级映射，提高 \code{__getitem__()} 性能；
\item 在样本读取阶段读取清晰 RGB 图像与对应深度缓存；
\item 将原始绝对坐标框转换为与图像空间对齐的 \codeword{xyxy} 表达，再统一转换为归一化 \codeword{xywh}；
\item 当指定 \code{img_size} 时，对图像和深度图共同执行 letterbox 或缩放处理，并同步变换检测框。
\end{enumerate}

值得强调的是，本文的输入空间处理保持了较强的一致性。无论是训练还是推理，系统都优先使用 letterbox，而不是简单拉伸。这是因为固定监控画面中横纵比往往固定，若直接拉伸至正方形，将改变车辆形状与车道透视比例，影响检测与天气判别的一致性。

\section{在线团雾合成方法}

为在缺乏大规模真实团雾标注数据的前提下训练雾分类与浓度回归任务，本文设计了位于 \texttt{src/model/fog\_augmentation.py} 中的在线造雾模块 \texttt{FogAugmentation}。该模块在 GPU 上实时运行，不预生成离线雾图，从而具有以下优势：一是节省磁盘空间；二是每轮训练都可动态生成不同参数组合，提升样本多样性；三是避免离线数据版本膨胀与管理复杂度。

在线造雾共构造三类天气样本：

\begin{itemize}[leftmargin=2em]
\item \texttt{clear}：无雾样本，保持原图像不变，并将 beta 标签置为 0；
\item \texttt{uniform}：均匀雾样本，整幅图像按统一散射参数退化；
\item \texttt{patchy}：团雾样本，在深度基础上叠加低频噪声，使局部区域雾强度发生变化。
\end{itemize}

在参数采样方面，系统根据 \texttt{Config} 中设置的物理范围对 beta 与大气光 A 进行随机采样，其中 \texttt{BETA\_MIN = 0.02}、\texttt{BETA\_MAX = 0.1}、\texttt{A\_MIN = 0.7}、\texttt{A\_MAX = 0.95}。对于 \texttt{uniform} 样本，系统根据深度图计算透射率，并对极端值进行裁剪，防止图像完全变白或几乎无雾。对于 \texttt{patchy} 样本，则先生成低分辨率随机噪声并双三次上采样到原图大小，再将该噪声作为有效深度调制因子，使浓雾和淡雾能够在同一画面局部并存。

从任务监督角度看，这一模块具有双重作用。首先，它提供雾类别标签，支撑分类头训练；其次，它输出实际用于生成雾图的 beta 值，支撑回归头学习。这意味着雾分类与 beta 回归不再依赖人工标注，而是由在线合成过程自动生成监督信号。虽然这种监督属于“合成标签”，与真实物理能见度并不严格等价，但它为模型学习“相对雾强度”提供了稳定、连续且可控的训练信号。

\section{在线造雾策略的合理性分析}

为何本文不直接收集雾天视频进行监督，而选择清晰数据加在线造雾？原因主要有三点。

首先，真实团雾样本获取成本高且分布稀疏。团雾往往发生在特定路段、特定季节和特定时段，不易形成稳定的大规模数据集。即使采集到视频，如何给每一帧准确赋予可比的浓度标签也是困难问题。

其次，在线造雾虽然不能完全还原真实世界，但能够用统一的物理启发生成可控监督，并通过大规模多样化组合提升模型学习到的退化规律。对于毕业设计与工程原型阶段，这是一种成本与效果之间较为合理的折中。

再次，本文关注的是“监测系统原型”而非高精度气象物理反演器。因此，beta 在本文中首先被视作“与雾强度正相关、可回归、可解释、可平滑”的相对指标，而不是需要绝对匹配气象站能见度的严格物理量。只要该指标能够稳定反映从清晰到严重退化的变化趋势，并能与视频检测和风险显示联动，就具备了实际价值。

\begin{figure}[htbp]
  \centering
  \includegraphics[width=0.95\textwidth]{fig_02_fog_augmentation_examples.png}
  \caption{在线造雾原理图与三类天气样例图}
  \label{fig:fog_augmentation_examples}
\end{figure}
